\section{The lower bound for independent transforms}\label{bch:proof:section}

A circuit can mix coordinates from all its input signals in a single
state. A low-rank correction to its output can do the same. To obtain
a bound proportional to the number of signals, we test the input
coordinates and Fourier outputs of each signal together. The rectangle
calculation in Section~\ref{geo:section} gives a lower bound on this
test even after a projection removes a small fraction of the dimensions.
We then show that a projection of the whole batch must preserve most
of these individual tests. Throughout the proof we work over the
algebraic closure $k$ of the coefficient field.

\subsection{A character table after projection}

For a matrix $U$ with columns $u_j$, its \emph{list tensor} is
$\mathcal D(U)=\sum_j u_j\otimes e_j\otimes e_j$. The label $j$ is
recorded twice so that we may place it at both ends of one side of a
rectangle. Put the vector leg at the center $O$ and the label legs on
the left side as in Figure~\ref{geo:rectangle-figure}. Write $V_-(U)$
for the spectral value $\Phi_S(\mathcal D(U))$ at these three sites
and $V_+(U)$ for its counterpart on the right. Set $V_\pm(0)=0$.
For a transform of order $N$, we use width one and height $L=3/\sqrt m$,
where $m=\log_2 N$. The triangle and rectangle have semiperimeters
\[
 \delta=\tfrac12\bigl(L+\sqrt{1+L^2}\bigr),
 \qquad \Delta=1+L .
\]
For the matrix $U=[I_N\ H]$ the list contains the coordinate vectors
and the columns of a character table. Its value on either triangle
is at most $2N$, the cost in spectral value of supplying the column
label itself. The next estimate retains nearly this value on a single
triangle after projection.

\begin{lemma}[Testing a projected character table]\label{bch:projected-list}
There is an absolute $N_0$ with the following property. Let $H$ be
the character table, in symmetric character coordinates, of a split
finite abelian group of order $N\ge N_0$ invertible in $k$. Put
$U=[I_N\ H]$ and $\varepsilon=2^{-100}$. Every linear map $Q$ on
the vector leg of rank at least $(1-\varepsilon)N$ satisfies
$V_-(QU)\ge1.8N$.
\end{lemma}

\begin{proof}
Choose a cyclic decomposition of the group and a primitive root on
each axis. The resulting symmetric bicharacter gives $H^{\mathsf T}=H$
and $H^{\mathsf T}H=NP$, where $P$ represents $g\mapsto-g$.
Put one copy of $\mathcal D(U)$ on each triangle and contract their
vector legs at $O$ by the usual bilinear pairing. The rectangle tensor
that remains is specified by the Gram matrix
\[
 U^{\mathsf T}U=\begin{pmatrix} I_N&H\\ H&NP\end{pmatrix}.
\]
The two diagonal blocks together have rectangle value $2N$, since
they supply a single common label up to local permutations and
nonzero scalings. The two off-diagonal blocks form a direct sum of
character-table tensors. Theorem~\ref{geo:saturation} gives each of
these price $2\delta m$ and hence rectangle value $2^{2\delta m/\Delta}$.
Subtracting the diagonal blocks from the Gram tensor and applying
\eqref{geo:sum} bounds the Gram value from below. On the other hand,
contracting at $O$ gives a price at most the sum of the two triangle
prices, by compatibility under insertion of trivial legs. Together
these estimates give
\[
   \bigl(V_-(U)V_+(U)\bigr)^{\delta/\Delta}
       \ge 2\,2^{2\delta m/\Delta}-2N .
\]
Put $a=m(2\delta-\Delta)/\Delta$. As each triangle value is at most
$2N$, the value on either one divided by $N$ is at least
$[2(1-2^{-a})]^{\Delta/\delta}/2$. The choice of height makes the
curvature term $a$ tend to $9/2$, while $\Delta/\delta\to2$.
The limiting lower bound is $2(1-2^{-9/2})^2>1.82$.

It remains to control the part removed by $Q$. Choose a linear map $E$
such that $EQ$ projects onto a complement to $\ker Q$ and put
$W=(I-EQ)U$. If its rank $r$ is positive, place a supplier halfway
from $O$ to the midpoint of the left side. A common label of size
$2N$ at the supplier and both left corners specifies a column of $W$.
Carry its coordinates in a basis of $\operatorname{im}W$ along a wire
of length $1/4$ and dimension $r$ to $O$. The triangle used by the
common label has semiperimeter
$\tau=L/2+\sqrt{1/16+L^2/4}$. Thus
\[
 \delta\log V_-(W)
 \le \tau\log(2N)+\tfrac14\log r
 \le \delta m+\tfrac92+\tau-25 .
\]
The last inequality uses $r\le2^{-100}N$ and
$\tau+1/4-\delta\le L^2/2=9/(2m)$. Since $\delta\to1/2$ and
$\tau\to1/4$, this gives $V_-(W)<0.01N$ for all sufficiently
large $N$. The estimate also holds when $r=0$. Now $U=EQU+W$,
so subadditivity and restriction at $O$ imply
$V_-(U)\le V_-(QU)+V_-(W)$.
\end{proof}

\subsection{Removing a rank correction from the whole batch}

Set $V=V_-$ and $p=L/\delta\in(0,1)$. We can annihilate a global
rank correction by projecting the output coordinates. This projection
may mix the blocks of independent signals. The next lemma extracts
one projection per block without increasing their total rank loss
or their total list value.

\begin{lemma}[A row quotient across independent lists]\label{bch:quotient-test}
Let $A=\bigoplus_i U_i$ have $h$ rows, in blocks of sizes $n_i$.
Suppose $A$ is in the closure of sums of a matrix of rank at most $r<h$
and a matrix with at most $T$ nonzero entries. There are $n_i$-square
matrices $Q_i$ such that
\[
 \sum_i(n_i-\operatorname{rank} Q_i)\le r,\qquad
 \sum_i V(Q_iU_i)\le h(T/h)^p .
\]
\end{lemma}

\begin{proof}
For an actual sum $B+P$, choose $h-r$ independent row vectors
annihilating $B$ and form the matrix $Q$ with these rows. Decompose
$QP$ into contributions from the rows of $P$. The contribution of
row $i$ has fixed vector leg, the $i$th column of $Q$, and two copies
of a label ranging over the $t_i$ nonzero entries of that row.
A wire of dimension $t_i$ and length $L$ between the two label
sites supplies this contribution, including its coefficients by a local
map. Its list value is at most $2^{L\log(t_i)/\delta}=t_i^p$ when
$t_i>0$, and zero when $t_i=0$. Subadditivity and concavity give
\[
 V(Q(B+P))=V(QP)
    \le\sum_i t_i^p\le h(T/h)^p .
\]

The annihilator can vary as $B+P$ approaches a border point. To see
that some quotient still gives the same bound there, consider the
Grassmannian of $(h-r)$-dimensional subspaces of $(k^h)^*$.
On each standard chart a basis matrix $Q$ depends regularly on the
subspace. For a matrix $M$, the condition $V(QM)\le h(T/h)^p$ is closed
by Theorem~\ref{geo:spectrum}, and changing the basis of the rows of
$Q$ preserves the list value. The chart conditions therefore define
a closed incidence set in the product of the Grassmannian with matrix
space. The Grassmannian is projective, so the image in matrix space
is closed. In particular the target $A$ admits such a $Q$.

At the block-diagonal target, put the rows of $Q$ in echelon form
and insert each into the row position of its pivot in a square matrix,
filling the other rows with zeros. This replacement preserves the list
value by local maps in both directions. The resulting matrix, still
denoted by $Q$, is upper triangular. Each pivot contributes one to
the rank of its diagonal block $Q_i=Q_{ii}$, so
$\sum_i\operatorname{rank}Q_i\ge h-r$.

The $(i,j)$ block of $QA$ is $Q_{ij}U_j$. Scale the vector block $i$
by $z^{-i}$ and the label block $j$ at one of its two sites by $z^j$.
These local maps are invertible for $z\ne0$. Each surviving block
acquires a factor $z^{j-i}$ with nonnegative exponent. At $z=0$ only
the diagonal blocks remain (Figure~\ref{bch:quotient-figure}). Closed
sublevels and additivity on direct sums give
$\sum_i V(Q_iU_i)\le V(QA)$.
\end{proof}

\begin{figure}[htb]
\centering
\begin{tikzpicture}[x=1.3cm,y=.65cm,font=\small]
  \foreach \off in {0,5} {
    \foreach \i in {0,1,2} {
      \fill[sea!12] (\off+\i,-\i) rectangle (\off+\i+1,-\i-1);
    }
    \draw[gray!65] (\off,0) rectangle (\off+3,-3);
    \foreach \i in {1,2}
      \draw[gray!65] (\off+\i,0)--(\off+\i,-3) (\off,-\i)--(\off+3,-\i);
    \foreach \i/\b in {0/1,1/2,2/3}
      \node at (\off+\i+.5,-\i-.5) {$Q_{\b}U_{\b}$};
  }
  \node at (1.5,.43) {Scaled quotient};
  \node at (6.5,.43) {Independent lists};
  \node at (1.5,-.5) {$zB_{12}$}; \node at (2.5,-.5) {$z^2B_{13}$};
  \node at (2.5,-1.5) {$zB_{23}$};
  \foreach \x/\y in {.5/-1.5,.5/-2.5,1.5/-2.5,5.5/-1.5,5.5/-2.5,6.5/-2.5,6.5/-.5,7.5/-.5,7.5/-1.5}
    \node[gray] at (\x,\y) {$0$};
  \draw[-{Stealth[length=2mm]},thick] (3.2,-1.5)--(4.8,-1.5)
    node[midway,above] {$z\to0$};
\end{tikzpicture}
\caption{An echelon quotient of three independent lists. The mixed
blocks $B_{ij}=Q_{ij}U_j$ vanish under scaling. Each pivot survives
in a shaded diagonal block, and the three surviving list values add.}
\label{bch:quotient-figure}
\end{figure}

\subsection{Repeating the transfer on a time coordinate}

The quotient test charges independent Fourier lists. To turn the number
of states into a rank budget for that test, we run the same transfer on
$d$ time coordinates. Making the time shift nilpotent changes
one column in a suitable basis. This costs at most one unit of rank
per state and truncates every dependency path at depth $d$, so each
row of the resulting transfer has small support.
For related time expansions of feedback circuits see
\cite[Sections~5--6]{FiniteFourier}.

\begin{lemma}[Time expansion]\label{bch:time-expansion}
Let $M\in\mathcal C_s(h)$ be the transfer of a binary linear system
and let $d\ge1$. Then $I_d\otimes M$ belongs to the closure of sums of
a matrix of rank at most $s$ and a matrix with at most $d2^{d-1}$
nonzero entries per row.
\end{lemma}

\begin{proof}
Write the state equations as $v=Cv+Ax$. Uniqueness at $x=0$ implies
that $I-C$ is invertible. Pad with zero states to obtain exactly $s$
states. Replace each state by $d$ copies, writing $C'=C\otimes I_d$
and $A'=A\otimes I_d$, and after each state apply a shift on its
$d$ time values. Consider two choices for this shift. First take
$R=I_d+D$, where $D$ sends $e_i$ to $e_{i+1}$ and kills $e_d$.
The vectors $e_1,Re_1,\ldots,R^{d-1}e_1$ form a basis in every
characteristic, by triangularity. In this basis take $R_0$ to act as
$R$ on the first $d-1$ vectors and to kill the last. Then $R_0^d=0$
and $\operatorname{rank}(R-R_0)=1$.

For $S=R$ or $R_0$, put $Z_S=I_s\otimes S$. The input-to-state
map $V_S$ solves $V_S=Z_S(C'V_S+A')$. For $R$, the equations are
triangular in the original time coordinates with invertible diagonal
blocks $I-C$. For $R_0$, $Z_{R_0}C'=C\otimes R_0$ is nilpotent.
Subtracting the two systems gives
\[
 (I-Z_R C')(V_R-V_{R_0})
       =(Z_R-Z_{R_0})(C'V_{R_0}+A').
\]
The right side has rank at most $s$, one per state. The same rank
bound holds after the output readout, since any direct readout of an
input is the same in both transfers.

In the transfer for $R_0$, any path visiting $d$ states has zero
contribution. Unroll a single readout into a binary dependency tree
truncated after $d-1$ state expansions along each path. At most
$2^{d-1}$ leaves can be original input coordinates. Each such coordinate
carries $d$ time values, so the row has at most $d2^{d-1}$ nonzeros.

In the original time basis the transfer for $R$ is lower triangular
in time with diagonal blocks $M$. Multiply each output time block $i$
by $t^{i-1}$ and each input time block $j$ by $t^{1-j}$. For $t\ne0$
this preserves rank and support. Sending $t$ to zero removes the
strictly lower blocks and leaves $I_d\otimes M$ in the claimed closure.
\end{proof}

\subsection{Comparing the two budgets}

\begin{proof}[Proof of Theorem~\ref{bch:uniform-robust}]
Let $h=JN$, $m=\log_2 N$, $d=\lceil\sqrt m/8\rceil$, and $h_d=dh$.
In symmetric character coordinates the inverse transform differs from
the forward transform by a permutation and a nonzero scalar. Such
permutations and scalings preserve all three budgets, so take each
block to be a forward transform $H_j$.

Suppose the target is in the geometric border of $M+B+Z$ as in
Theorem~\ref{bch:uniform-robust}. For an actual sum, apply
Lemma~\ref{bch:time-expansion} to $M$ and repeat $B,Z$ on the $d$
time coordinates. Transpose the result and adjoin $I_{h_d}$ as
input columns. Transposition preserves the total number of nonzero
entries, and the identity adds $h_d$ entries. This operation on
$M+B+Z$ is affine in the sum, so it maps the joint border into
the closure of rank-plus-sparse sums with the same budgets. At the
target the resulting list is, up to permutations, the direct sum of
$d$ copies of each $[I_N\ H_j]$. The rank and total support budgets are
\[
 r=s+d\lfloor\beta h\rfloor,\qquad
 T=h_d(1+d2^{d-1})+d\lfloor\eta h\,2^{\sqrt m/8}\rfloor .
\]

Take $\beta=\varepsilon/400$ and $\eta\le1$. If
$s\le\varepsilon h_d/400$, the sum of rank deficits in
Lemma~\ref{bch:quotient-test} is at most $\varepsilon h_d/200$.
At least a fraction $199/200$ of the $Jd$ blocks therefore retain
rank at least $(1-\varepsilon)N$. For all sufficiently large $N$,
Lemma~\ref{bch:projected-list} gives total list value at least
$1.8(199/200)h_d>1.79h_d$. But Lemma~\ref{bch:quotient-test}
bounds it by $h_d(T/h_d)^p$, where
\[
 (T/h_d)^p\le (1+d2^{d-1}+2^{\sqrt m/8})^p
       \longrightarrow 2^{3/4}<1.79 .
\]
Indeed $p\sqrt m\to6$ and $d/\sqrt m\to1/8$. This contradiction
gives $s>\varepsilon h\sqrt m/3200$ uniformly in $J$ and the field.

For the finitely many smaller orders, we find two disjoint sets of
input columns on which many output rows still have full rank. In each
group choose a cyclic factor of order $e\ge2$ and let
$\ell=\lfloor e/2\rfloor$. The first $\ell$ exponent rows paired
with either of two disjoint sets of $\ell$ nodes give invertible
Vandermonde minors. Tensoring with the other cyclic factors and taking
the direct sum over the groups gives a common set of $u\ge h/3$
rows with full row rank on each of two disjoint column sets.

For any circuit readout, the direct-input part of each selected row
misses at least one of the two column sets. A state readout has zero
direct-input part and misses both. On at least $u/2$ rows one common
column set is missed. The transfer restricted to these rows and that
column set has rank at most $s$ by factorization through the states.
The corrections add at most $\beta h+|\operatorname{supp} Z|$ to
the rank. There are finitely many choices of rows and column set at
each size, so the union of these rank conditions is closed. Full row
rank at the target now gives $u/2\le s+\beta h+\eta h2^{\sqrt m/8}$.
Shrink $\eta$ once so that $\beta+\eta 2^{\sqrt m/8}\le1/12$ for
every order in this finite range. Then $s\ge h/12$. Decreasing the
absolute constant in the lower bound proves the assertion at all orders.
An acyclic circuit of size $s$ is a uniquely solvable system with $s$
states. Taking $B=Z=0$ gives the lower bound of Theorem~\ref{thm:main}.
\end{proof}
